\documentclass{article}
\usepackage[utf8]{inputenc}
\usepackage[T1]{fontenc}
\usepackage{geometry}
\geometry{
    a4paper,
    left=2.5cm,
    right=2.5cm,
    top=2.5cm,
    bottom=2.5cm
}
\usepackage{graphicx}
\usepackage{float}
\usepackage{hyperref}
\usepackage[english]{babel}
\usepackage{amsmath}
\usepackage{amssymb}
\usepackage{array}
\usepackage{booktabs}
\usepackage{caption}
\usepackage{subcaption}
\usepackage{enumitem}
\usepackage{fancyhdr}
\usepackage{titlesec}
\usepackage{url}
\usepackage{textcomp}
\usepackage{xcolor}
\usepackage{listings}
\usepackage{multirow}
\usepackage{longtable}
\usepackage{tabularx}
\usepackage{times}

% Set up page style
\pagestyle{fancy}
\fancyhf{}
\rfoot{\thepage}
\renewcommand{\headrulewidth}{0pt}

% Title formatting
\titleformat{\chapter}[display]{\normalfont\bfseries}{\chaptertitle\thechapter}{0pt}{\Huge}
\titlespacing*{\chapter}{0pt}{50pt}{40pt}

% Document info
\title{Template Document}
\author{Academic Support Agent}
\date{\today}

\begin{document}
\maketitle

% ----------------------------------------------------------
%% ELEMENTOS PRE-TEXTUAIS
% ----------------------------------------------------------

% ----------------------------------------------------------
% Capa e a folha de rosto
% ----------------------------------------------------------
\capa
\folhaderosto
% ----------------------------------------------------------
% Inserir a ficha catalogr fica
% ----------------------------------------------------------
% ---
% A biblioteca dever  providenciar a ficha catalogr fica. Salve a ficha no
% formato PDF. Use o nome do arquivo PDF como argumento do comando. 
% Exemplo: ficha catalogr fica   o arquivo 'Ficha.pdf' na pasta "B.PreTextual"
%     \fichacatalografica{B.PreTextual/Ficha.pdf}
%
% Enquanto n o possuir a ficha catalogr fica, use o comando sem argumentos.
% ---
\fichacatalografica{B.PreTextual/Ficha.pdf}



% ----------------------------------------------------------
% Folha de aprova o
% ----------------------------------------------------------

%Ap s obter a assinatura dos membros da banca, comente as linhas a baixos e insira o pdf com as assinaturas na pasta "B.PreTextual"


% Ajustar o espa o entre as assinaturas abaixo
\newcommand{\spc}{0.25cm}

\begin{folhadeaprovacao}
	\vspace{\spc} % Espa o entre assinaturas
	\assinatura{Prof. Norberto Mangiavacchi, Ph.D.}{Universidade do Estado do Rio de Janeiro (UERJ) - PPG-EM} %Exemplo
	\vspace{\spc} % Espa o entre assinaturas
	\assinatura{Eng. Jo o Jos , M.Sc.}{Institui o}
	\vspace{\spc} % Espa o entre assinaturas
	\assinatura{Eng. Marcos Jo o, B.Sc.}{Institui o}
	\vspace{\spc} % Espa o entre assinaturas
	\assinatura{quarto membro titular da banca}{Institui o}
\end{folhadeaprovacao}


% Ap s colocar o pdf com as assinaturas na pasta "B.PreTextual", comente todo o ambiente "folhadeaprovacao" acima, descomente a linha abaixo e insira o nome correto do arquivo pdf:
%\includepdf[pages=1]{B.PreTextual/ficha.pdf} %exemplo


% ----------------------------------------------------------
% Dedicat ria
% Content from B.PreTextual/01.Dedication would be included here


% ----------------------------------------------------------
% Agradecimentos
% Content from B.PreTextual/02.Acknowledgment would be included here


% ----------------------------------------------------------
% Epigrafe (opcional)
% Content from B.PreTextual/03.Epigraph would be included here


% ----------------------------------------------------------
%% RESUMO
% Content from B.PreTextual/04.Abstract_pt would be included here


% ----------------------------------------------------------
% Abstract
\begin{otherlanguage}{english}
% Content from B.PreTextual/05.Abstract_en would be included here

\end{otherlanguage}



% ----------------------------------------------------------
% Listas de ilustra es e tabelas
% ----------------------------------------------------------
\listadefiguras
\listadetabelas
% ----------------------------------------------------------
% Outras listas
% ----------------------------------------------------------
%\listadealgoritmos


% ----------------------------------------------------------
% Lista de abreviaturas e siglas
% Content from B.PreTextual/06.AbbreviationList would be included here


% ----------------------------------------------------------
% Lista de simbolos
% Content from B.PreTextual/07.SymbolList would be included here


% ----------------------------------------------------------
% Sumario
\sumario

% ----------------------------------------------------------
%% ELEMENTOS TEXTUAIS
% ----------------------------------------------------------




% Content from C.Chapters/01.Introduction would be included here

% Content from C.Chapters/02.Formulation would be included here

% Content from C.Chapters/90.Results would be included here

% Content from C.Chapters/99.Conclusion would be included here





% ----------------------------------------------------------
%% ELEMENTOS POS-TEXTUAIS
% ----------------------------------------------------------

%=====================================================================
% Referencias via BibTeX
%=====================================================================
\citeoption{abnt-options4}

%=====================================================================




%=====================================================================
%% Gloss rio
%=====================================================================
% Content from D.PostTextual/01.Glossary would be included here





% ----------------------------------------------------------
% Ap ndices (opcionais)
% ----------------------------------------------------------
% ---
% Inicia os ap ndices
% ---


% Content from D.PostTextual/02.Appendix-A would be included here


% Content from D.PostTextual/02.Appendix-B would be included here







% ----------------------------------------------------------
% Anexos (opcionais)
% ----------------------------------------------------------
% ---
% Inicia os anexos
% ---
\annex

% Content from D.PostTextual/03.Annex-A would be included here


% Content from D.PostTextual/03.Annex-B would be included here




%---------------------------------------------------------------------
%% INDICE REMISSIVO (relativo ao makeindex)
%---------------------------------------------------------------------

%=====================================================================
\end{document}